\subsection{The D6 period and the prepotential constant}
\label{subsec:x9-d6-period}

The symplectic pairing and the reality of the marked continuation fix the
eighth period. Curvature and flux shifts are essential in making this
an integral statement. We first specify the D-brane basis; the D6 label
below denotes a six-brane wrapping the threefold, not the toric divisor
$D_6$.

\paragraph{The integral symplectic frame.}
Appendix~\ref{app:integral-charges} constructs the primitive
topological charge generators $\mathsf m_0,\mathsf m_a,\mathsf e_0,
\mathsf e_a$, with pairing minus the Euler form. Use the ordered
divisor bases $L_a=H_a$ on the resolution and $L_a=K_a$ on the
smooth side, and write $c_a=c_2(X)L_a$, $\kappa_{abc}=L_aL_bL_c$.

Define the integers
\begin{equation}
 I_a=\frac{\kappa_{aaa}}6+\frac{c_a}{12},\qquad
 h_{ab}=\frac{\kappa_{abb}-\kappa_{aab}}2.
 \label{eq:x9-d6-rr-integers}
\end{equation}
Riemann--Roch gives $I_a\in\ZZ$ and $h_{ab}\in\ZZ$.
The unimodular correction
\begin{equation}
 \mathsf m'_0=\mathsf m_0,\qquad
 \mathsf m'_a=\mathsf m_a+I_a\mathsf e_0+\sum_{b>a}h_{ab}\mathsf e_b
 \label{eq:x9-d6-integral-correction}
\end{equation}
puts the pairing in standard symplectic form, leaving the electric
generators fixed \cite[Equations~(2.20)--(2.24)]{Demirtas:2024Mirror}.
The two index rows are $(2,3,9,26,36,50)$ and $(3,9,26)$.
In this convention the prepotential and normalized periods are
\begin{equation}
 \begin{aligned}
 F^X(t)&=-\frac16\kappa_{abc}t_at_bt_c
       +\frac12 a_{ab}t_at_b+\frac1{24}c_at_a
       +\frac{\chi(X)\zeta(3)}{2(2\pi i)^3}
       +F^X_{\rm inst}(t),\\
 G^X&=2F^X-t_aF^X_a,\qquad F^X_a=\partial_{t_a}F^X,\qquad
 a_{ab}=\begin{cases}
 \kappa_{aab}/2,&a\geq b,\\
 \kappa_{abb}/2,&a<b.
 \end{cases}
 \end{aligned}
 \label{eq:x9-d6-full-prepotential}
\end{equation}
Repeated indices are summed except in the definition of $a_{ab}$.
The period $G^X$ belongs to $\mathsf m_0$, which was not changed by
\eqref{eq:x9-d6-integral-correction}. Equivalently, in the
\emph{unsubtracted} Frobenius derivatives,
\begin{equation}
 W_XG^X=\frac16\kappa_{abc}\varpi^{abc}_X
                 +\frac1{12}c_a\varpi^a_X,
 \label{eq:x9-d6-frobenius-normalization}
\end{equation}
where $\varpi^{abc}_X$ includes the factors $(2\pi i)^{-3}$
\cite[Equations~(2.29) and~(4.19)]{Demirtas:2024Mirror}.
Removing the $\zeta(3)$ term would change this marked period.

The surviving block of $a_{ab}$ on the resolution equals the
smooth-side matrix, just as the cubic and $c_2$ rows agree:
\begin{equation}
 (a^{\rm res}_{i+1,j+1})_{i,j=1}^3=a^{\rm sm}
 =\begin{pmatrix}
 0&3&17/2\\3&7&49/2\\17/2&49/2&83/2
 \end{pmatrix}.
 \label{eq:x9-d6-quadratic-matrix}
\end{equation}
Proposition~\ref{prop:x9-neutral-period-descent} therefore identifies
the seven unnormalized integral periods belonging to
$\mathsf e_0,\mathsf e_2,\mathsf e_3,\mathsf e_4,
\mathsf m'_2,\mathsf m'_3,\mathsf m'_4$ with their smooth-side counterparts.
For example their magnetic components are
\begin{equation}
 W F_i=-\Pi^{\rm mag}_i+\frac{c_i}{24}W
                         +\sum_j a_{ij}\Pi^{\rm el}_j.
 \label{eq:x9-d6-seven-integral-periods}
\end{equation}
Local electric terms on the resolution vanish on $\Sigma$.

\begin{proposition}\label{prop:x9-d6-period}
On the marked small-bulk branch of $\Sigma$, the reduced integral
symplectic period vector agrees with that of $X_{9,t}$. In the
conventions \eqref{eq:x9-d6-full-prepotential},
\begin{equation}
 G^{\rm res}|_\Sigma=G^{\rm sm},\qquad
 F^{\rm res}|_\Sigma=F^{\rm sm},\qquad
 F^{\rm res}_{\rm inst}|_\Sigma-F^{\rm sm}_{\rm inst}
       =-\frac{4\zeta(3)}{(2\pi i)^3}.
 \label{eq:x9-d6-period-and-constant}
\end{equation}
The integral map sends the surviving classes
$(\mathsf m'_0,\mathsf m'_2,\mathsf m'_3,\mathsf m'_4;
\mathsf e_0,\mathsf e_2,\mathsf e_3,\mathsf e_4)$ to the correspondingly
ordered smooth-side basis, in $V^\perp/V$.
\end{proposition}

\begin{proof}
Conifold excision identifies $H_3$ of the complement of the vanishing
spheres with $V^\perp$. Filling by the small-resolution neighborhoods
quotients it by the images of the boundary three-spheres, namely $V$.
Since $V$ is primitive, this is the free integral symplectic quotient
$V^\perp/V$. The residue-preserving birational identification in
Proposition~\ref{prop:x9-mirror-mutation} preserves this pairing.
The corrected basis \eqref{eq:x9-d6-integral-correction} makes the
quotient explicit: discard $\mathsf m'_1,\mathsf m'_5,\mathsf m'_6$ and quotient by
$\mathsf e_1,\mathsf e_5,\mathsf e_6$. The eight classes stated in the proposition give a
unimodular symplectic basis.

The D6 class $\mathsf m_0$ lies in $V^\perp$, since every vanishing charge
is pure D2. Its period thus descends. The difference of its descended
class and the smooth-side $\mathsf m_0$ pairs trivially with each of the seven
already identified classes. Their common annihilator is the line
spanned by $\mathsf e_0$. Consequently
\begin{equation}
 W\bigl(G^{\rm res}|_\Sigma-G^{\rm sm}\bigr)=cW,
 \qquad c\in\ZZ.
 \label{eq:x9-d6-scalar-ambiguity}
\end{equation}
This step uses the polarization, not just the dimension of the period
space. A real coefficient would already suffice below.

We determine $c$ on the positive real branch. Subtract the quadratic
term from the prepotential and write
$\widetilde F_a=F_a-a_{ab}t_b$; this rational symplectic change
does not alter $G$. The electric and second-derivative Abel
continuation in Appendix~\ref{app:mirror-periods}, along
$z_1=z_5=z_6=\rho\uparrow1$ with fixed sufficiently small positive
$(v,a,b)$, has
\begin{equation}
 W\in\mathbb R\setminus\{0\},\qquad t_a\in i\mathbb R,\qquad
 \widetilde F_a\in\mathbb R.
 \label{eq:x9-d6-real-continuation}
\end{equation}
The Gamma coefficients and logarithms are real before their explicit
powers of $2\pi i$ are inserted. The convergent weighted sums fix
these reality properties on the continuation, rather than selecting
a new branch at the endpoint. The initial large-volume value of $G$
is purely imaginary by \eqref{eq:x9-d6-full-prepotential}; the
quadratic term cancels from $2F-t_aF_a$. Special geometry gives
\begin{equation}
 dG=\sum_a\bigl(\widetilde F_a\,dt_a-t_a\,d\widetilde F_a\bigr).
 \label{eq:x9-d6-reality-differential}
\end{equation}
Its right-hand side is purely imaginary along this path. Hence the
continued $G^{\rm res}$ is purely imaginary, including its finite
endpoint supplied by conifold descent. This argument requires no
termwise limit of an uncontracted third-derivative series. The
smooth-side $G^{\rm sm}$ is likewise purely imaginary on its
positive small-parameter branch. Equation~\eqref{eq:x9-d6-scalar-ambiguity}
then forces the real integer $c$ to vanish. Analytic continuation
from this open real locus gives the identity throughout the marked
complex neighborhood.

The seven-period result, together with
\eqref{eq:x9-d6-quadratic-matrix}, makes all three derivatives of
$F^{\rm res}|_\Sigma-F^{\rm sm}$ zero. Equality of $G$ fixes this
constant to zero. Finally
$\chi(\wideX_9)=-232$ and $\chi(X_{9,t})=-240$; subtracting the
perturbative terms in \eqref{eq:x9-d6-full-prepotential} gives the
last equality in \eqref{eq:x9-d6-period-and-constant}.
\end{proof}

\paragraph{Scalar and electromagnetic couplings.}
The eighth-period match determines the two-derivative vector action,
including its coupling to the graviphoton. In the common normalization,
\begin{equation}
 e^{-K_{\rm vec}}=i\bigl(\overline{X^A}F_A-X^A\overline{F_A}\bigr),
 \qquad (X^A)=(1,\tau_i),\quad (F_A)=(G,F_i),
 \label{eq:x9-d6-special-kahler}
\end{equation}
has the same restriction on both sides. Consequently the scalar kinetic
metric $g_{i\bar j}=\partial_{\tau_i}\partial_{\bar\tau_j}K_{\rm vec}$
agrees. The terms involving the three vanishing electric coordinates
drop out of the original pairing.

Set $\mathscr F(X)=(X^0)^2F(X^i/X^0)$ for the homogeneous prepotential, and write
$\mathscr F_{AB}=\partial_A\partial_B\mathscr F$ for $A,B=0,\ldots,3$.
The supergravity period matrix is
\begin{equation}
 \mathcal N_{AB}=\overline{\mathscr F}_{AB}
 +2i\frac{(\operatorname{Im}\mathscr F_{AC})X^C
                (\operatorname{Im}\mathscr F_{BD})X^D}
               {X^E(\operatorname{Im}\mathscr F_{EF})X^F}.
 \label{eq:x9-electromagnetic-couplings}
\end{equation}
In the convention where $-\operatorname{Im}\mathcal N$ is the
positive gauge-kinetic form, $\operatorname{Re}\mathcal N$ gives
the theta couplings \cite[Equations~(1.1) and~(4.22)]{crapsetal1997specialgeometry}.
Equality of the reduced prepotentials therefore identifies both matrices,
including graviphoton mixing, in the stated integral electric frame.
These are the surviving couplings after conifold reduction, not the
divergent charged block in \eqref{eq:x9-threshold-matrix}.

The eighth period matters because a constant in the affine $F$ becomes
a quadratic term in the homogeneous $\mathscr F$. Its imaginary part
can change $K_{\rm vec}$ and $\mathcal N$, even though it leaves the
affine Yukawa couplings unchanged. Real quadratic changes also affect
the electromagnetic frame. The integral marking and constant matching
remove these ambiguities. This determines the genus-zero vector-sector
couplings on the marked branch, not a hypermultiplet metric or a
higher-derivative effective action.
